% Extensión estimada: algo más de media página, tres párrafos.
% Fuentes: Codigo/CLAUDE.md (GPU de cada clúster y del equipo local, bf16 en Ampere o posterior);
% Codigo/requirements.txt y src/models/{classifiers,dermlip,mini_vlm}.py (bibliotecas); scripts/measure_cost.py
% (costo de los clasificadores en la RTX 4070, lotes de 32); Apéndice E (ap:reproducibilidad).
% Revisión del tutor del 21/10/2026: sección nueva que reúne el hardware, las semillas y las bibliotecas, antes
% dispersos en 3.1, 3.3, 3.6 y 3.8. Sin versiones de CUDA ni herramientas internas (regla R1 del autor).
\section{Entorno experimental y recursos computacionales}
\label{sec:met_entorno}

Los experimentos se ejecutaron en tres GPU, todas dentro de la restricción de la Sección~\ref{sec:met_tareas}. Los
clasificadores supervisados se entrenaron en una RTX 4080 de 16~GB, en el clúster TOKO de la Universidad Nacional de
Cuyo. El ajuste de los \glspl{vlm} y la generación de los reportes se ejecutaron en una A30 de 24~GB, en el clúster
Mendieta de la Universidad Nacional de Córdoba. Por último, una RTX 4070 de 12~GB, en un equipo personal, se usó para
calcular las métricas de generación y para medir el costo de los clasificadores. Las tres GPU tienen arquitectura Ampere
o posterior, lo que permite entrenar en precisión \emph{bfloat16} (bf16) en todos los casos.

En cuanto a la medición del costo (Sección~\ref{sec:met_costo}), el tiempo de inferencia y la cantidad de imágenes por
segundo de los clasificadores se midieron en la RTX 4070, con lotes de 32 imágenes. El tiempo por reporte de los
generadores, en cambio, se tomó de las ejecuciones de la generación en la A30. Como las dos mediciones se hicieron en
GPU distintas, la comparación entre clasificadores y generadores es solo de orden de magnitud.

El código usa PyTorch y torchvision para los clasificadores supervisados, las bibliotecas \texttt{transformers} y
\texttt{peft} de Hugging Face para los \glspl{vlm} generativos y LoRA, \texttt{open\_clip} para DermLIP y scikit-learn
para las cabezas clásicas, la partición y las métricas. Todos los entrenamientos fijan la semilla 42, salvo las
réplicas con las semillas 43 y 44 de la Sección~\ref{sec:met_supervisados_ensamble}. El
Apéndice~\ref{ap:reproducibilidad} reúne lo necesario para repetir cada experimento.
